% Pursuing Stacks (cotes 134-2 to 134-8): archive cote 134-8, site page 18 (PDF page 19), typescript page 572.
% typescript. Draft compared: ps7.tex:825-887.
% AI-transcribed by Claude Opus 5.5 (claude-opus-5-5, adaptive thinking, effort high), Message Batches API, batch ps4 (msgbatch_012Aoc1NuskATPsyBeKSYK7J), 2026-11-12; not yet checked by eye.
\scanpage{134-8/18}
\mspage{572}

These reflections make it quite clear that the \xout{relevant} notions
of properness and\inserted{of} smoothness for maps in (Cat), relevant for \deleted{working
out} a duality formalism, have still to be worked out. Two basic requi-%
rements to be kept in mind are the following: 1) for a proper map \xout{fxx}
f: A $\longrightarrow$ B, and any ring of coefficients k, the functor
\[ Rf_{*} : D^{+}(A^{\wedge}_{k}) \longrightarrow D^{+}(B^{\wedge}_{k}) \]
should admit a right adjoint $f^{!}$ , and 2) for a smooth map f factored
as $f = g\circ i$ , with g proper and i an "open immersion" , the compositi-%
% NOTE: the end of this line runs off the right edge of the scan; one character after "compositi" is cut
on \deleted{f}\inserted{$f^{!}\overset{\mathrm{def}}{=}$}$i^{*}\circ g^{!}$\hand{$\leftarrow$} should be \xout{isomorphic to} expressible as
\margin{$f^{!}\overset{\mathrm{def}}{=} i^{*}\circ g^{!}$}
\[ (5\unsure{)} \qquad f^{!} : \quad K^{\bullet} \mapsto f^{*}(K)\overset{L}{\otimes}_{k} T_{f}(k)\hand{$[$}d_{f}\hand{$]$} , \]
where $T_{f}(k)$ is a presheaf of k-modules on A locally isomorphic to
the constant \xout{sheaf} presheaf $k_{A}$ ($T_{f}(k)$ may be called the \uline{orientation}
\uline{sheaf} for f, with coefficients in k), and $d_{f}$\hand{|}is a\xout{n} natural integer
(which may be called the \uline{relative dimension} of A over B, or of f).
\hand{(}For simplicity I assume in 2) that A is connected, otherwise $d_{f}$
should be viewed as a function on the set of connected components of
A .\hand{)} This again should give the correct relationship, for f \xout{\illegible{1}}\inserted[typed]{as above}\hand{,}
between the (for the time being hypothetical) $Rf_{!}$ ($\overset{\mathrm{def}}{=} Rg_{*}\circ Li_{!}$ ) and
\xout{the} our anodyne $Lf_{!}$ , for an argument L. in $D^{\over{\bullet}{b}}(A^{\wedge}_{k})$ say :
\[ (6) \qquad Rf_{!}(L.) \simeq Lf_{!}(L.\otimes T_{f}^{-1})\hand{$[$}-d_{f}\hand{$]$} , \]
where the \xout{first} left hand side is just $Rf_{*}(L.)$, if we assume moreover
f to be proper.

This precise relationship between the two possible versions of
an\xout{f} $f_{!}$ operation \xout{atxxx} between derived categories, namely $Lf_{!}$
embodying homology, defined for any map f in (Cat), and $Rf_{!}$ embodying
"cohomology with proper supports" , defined for a map which may be
factored as $g\circ i$ with g\hand{/}"proper"and i an (open, if we like) immersion,
relation valid if f is moreover assumed to be"smooth", greatly relaxes
\xout{the}\inserted{yesterday's} perplexity, coming from a\inserted{partial} confusion in my mind between the operati-%
ons $Rf_{!}$ and $Lf_{!}$ . (\deleted{where}\inserted{Beware} the notation $Rf_{!}$ \deleted{being}\inserted{is} an abuse,\xout{\illegible{1}} as
\inserted[typed]{\xout{in the present context}}\inserted{here,} it does'nt mean at all anything like the \xout{derived} right derived functor
of the functor $f_{!}$ \hand{!}) At the same time, I feel a lot less dubious now
about the existence of a \xout{duality} "six operations" duality formalism in
the (Cat) context - I am pretty much convinced now that such formalism
exists indeed. The main specific work ahead is to get hold of the
relevant notions of proper and smooth maps. THe demands we have on

\margin{18}
